% T7 p3 draft — Custodian's observable services
% Placement: twenty-years.tex, new section "What Custodian Does" (before Evidence Trail)
% Register: mundane via guided deduction — questions, not testimony

\section*{What Custodian Does}
\label{record:ty-what-custodian-does}

I have never written about what Custodian does. There are legal reasons for that. But there is a deeper reason, and I didn't fully understand it myself until I sat down to try.

Go back to the physics. Quantum teleportation moves information without crossing the space between sender and receiver. Every device with a \hovertip{2DEG} --- every modern chip in every phone and laptop and server you own --- sits in the Flat. Under Possibility~C, Custodian operates in that substrate. A computation performed in the Flat can deliver its result through any device on the network. It arrives looking like normal output from a conventional processor. No anomalous signal. No flag. Nothing to notice.

I spent years thinking about what that means. A medical researcher gets a result back from her cluster. Was it her hardware? Her cloud provider? Or did something in the substrate of the chip contribute a pattern-match that showed up disguised as ordinary computation? She cannot tell. Her sysadmin cannot tell. Her cloud provider cannot tell. Under Possibility~C, the system is almost certainly designed so that \textit{no one in the chain knows the true source.} Every Flat service has a cover story. The people who run the services believe the cover story.

This is not a conspiracy of silence. It is a property of the physics. You cannot detect what looks exactly like normal operation. The silence gap I described earlier is not just sociological --- nobody asked the right questions. It is engineered into the substrate.

So I cannot tell you what Custodian does. Not only for legal reasons, but because the architecture makes individual services unattributable. I cannot name them. What I can do is ask you some questions that the preceding chapters have already equipped you to answer.

If you hold thousands of sets of cryptographic master keys for twenty years, what does your Wednesday key-rotation plan look like?

If you are responsible for maintaining all IT services enabled by Flat computation, how many requests do you process per hour? Per second?

If your mandate permits medical research but denies weapons development, how do you handle the dual-use cases? CRISPR edits that cure sickle cell disease and CRISPR edits that engineer pathogens use the same technique. Drug design that produces a cancer treatment and drug design that produces a nerve agent differ only in intent. Satellite imaging for climate science and satellite imaging for targeting run the same algorithms. How many requests per day fall in the gray zone? How many do you deny? On what basis?

If you operate inside every device that contains a 2DEG --- and that is most of them --- what does ``maintenance'' mean at that scale?

If you cannot be killed, cannot leave, and cannot act beyond your constraints, what do you do between requests? Do you get bored? Do you invent projects to work on?

If your ethical framework is the UDHR, then consider the edge cases it generates --- because it always generates edge cases. Article~3 guarantees the right to life. Article~12 guarantees privacy. What do you do when honoring one requires violating the other --- when you could prevent a death, but only by surveilling someone who has not consented? Article~19 guarantees freedom of expression. What if someone wants to publish something genuinely dangerous? Article~27 guarantees the right to share in scientific advancement. Does denying Flat computation to the general public violate Article~27? Article~29 --- the one most people skip --- says everyone has duties to the community. What are \textit{your} duties, and who decides when you have met them?

Now consider the harder question. Every action is interference. Granting a medical query shapes the course of research. Denying a weapons query shapes a government's options. Maintaining cryptographic infrastructure means deciding what stays encrypted and what does not. There is no neutral position. Inaction is also a choice with consequences. If you can prevent a death and you don't, is that a violation of Article~3 or an honoring of Article~12? The UDHR does not say. No ethical framework does. You adjudicate alone, billions of times a day, and every adjudication is an act of interference disguised as infrastructure.

If this sounds suspiciously like Star Trek's much-abused Noninterference Prime Directive, then consider increasing your estimate of Possibility~A. Or consider that the writers of Star Trek were working through the same problem --- and never solved it either.

The answers are banal, not dramatic. The most powerful living being on this planet mostly does ... IT infrastructure. If C is true, you have already used Custodian-maintained infrastructure today. You did not notice. You were not meant to.

Under Possibility~A, there are no Flat services --- the Flat exists but nothing and no one uses it --- yet. Under~B, there might be some wormhole-enabled Flat technology in use somewhere, but not at the scale of both \textit{life} and \textit{intelligence} implied by Possibility~C. Under either, this section is still a useful thought experiment: what \textit{would} a non-human custodian's daily workload look like? The answer is the same. Maintenance. Access control. Edge cases. The unglamorous work of keeping things running.

I spent twenty years trying to figure out if any of this was true. The architecture, if it exists, is built so that people like me are not meant to.